%
% teil1.tex -- Herleitung der Bessel-Differentialgleichung
%
% (c) 2026 Gian Cavegn, OST Ostschweizer Fachhochschule
%
% !TEX root = ../../buch.tex
% !TEX encoding = UTF-8
%
\section{Herleitung der Bessel-Differentialgleichung
\label{bessel:section:herleitung}}
\kopfrechts{Herleitung}
% Funktion dieses Abschnitts:
% Zeigen, wie die Bessel'sche DGL aus einem konkreten physikalischen
% Problem entsteht. Das ist der "Aha"-Moment des Kapitels.
% Roter Faden: Wellengleichung -> stationäre Lösung (Helmholtz) ->
% Zylinderkoordinaten -> Separation -> Bessel-DGL für den Radialteil.

Bevor wir die Bessel'sche Differentialgleichung analysieren, soll gezeigt werden, woher sie kommt.
Wir leiten sie aus einem klassischen physikalischen Problem her: der Helmholtz-Gleichung in Zylinderkoordinaten.

\subsection{Die Wellengleichung und die Helmholtz-Gleichung
\label{bessel:subsection:helmholtz}}
% Inhalt:
% - Wellengleichung in 2D/3D: u_tt = c^2 \Delta u.
% - Stationärer Ansatz mit harmonischer Zeitabhängigkeit:
%   u(\vec x, t) = U(\vec x) e^{i\omega t}
% - Einsetzen liefert die Helmholtz-Gleichung \Delta U + k^2 U = 0
%   mit k = \omega / c.

Die räumliche Ausbreitung von Wellen wird durch die Wellengleichung
\begin{equation}
\frac{\partial^2 u}{\partial t^2}
=
c^2 \Delta u
\label{bessel:equation:wellen}
\end{equation}
beschrieben, wobei $\Delta$ den Laplace-Operator und $c$ die Ausbreitungsgeschwindigkeit bezeichnet.
Wir suchen Lösungen mit harmonischer Zeitabhängigkeit der Form
\begin{equation}
u(\vec x, t)
=
U(\vec x)\, e^{i\omega t}.
\label{bessel:equation:ansatz-zeit}
\end{equation}
Einsetzen in~\eqref{bessel:equation:wellen} liefert für die räumliche Funktion $U$ die Helmholtz-Gleichung
\begin{equation}
\Delta U + k^2 U = 0,
\qquad
k = \frac{\omega}{c}.
\label{bessel:equation:helmholtz}
\end{equation}
% [TODO: kurzer Kommentar dazu, dass k die Wellenzahl ist]

\subsection{Zylinderkoordinaten
\label{bessel:subsection:zylinderkoordinaten}}
% Inhalt:
% - Laplace-Operator in Zylinderkoordinaten (r, phi, z) explizit hinschreiben.
% - Argumentieren, warum diese Koordinaten bei rotationssymmetrischen
%   Problemen die "natürlichen" sind.

Bei Problemen mit zylindrischer Symmetrie ist es zweckmässig, in Zylinderkoordinaten $(r, \varphi, z)$ zu rechnen.
Der Laplace-Operator nimmt in diesen Koordinaten die Form
\begin{equation}
\Delta U
=
\frac{1}{r}\frac{\partial}{\partial r}\!\left(r\frac{\partial U}{\partial r}\right)
+ \frac{1}{r^2}\frac{\partial^2 U}{\partial \varphi^2}
+ \frac{\partial^2 U}{\partial z^2}
\label{bessel:equation:laplace-zylinder}
\end{equation}
an.
% [TODO: kurze Herleitung oder Verweis auf Standardliteratur]

\subsection{Separationsansatz
\label{bessel:subsection:separation}}
% Inhalt:
% - Produktansatz U(r,\varphi,z) = R(r) \Phi(\varphi) Z(z)
% - In Helmholtz einsetzen, durch U teilen.
% - Drei einzelne ODEs herausarbeiten, mit Separationskonstanten.
% - Die phi-Gleichung -> sin/cos mit ganzzahliger Periodizität,
%   ergibt Konstante n^2 (oder allgemeiner \nu^2).
% - Die z-Gleichung -> Exponential/Trigonometrisch.
% - Übrig bleibt die Radialgleichung für R(r).

Wir machen den Produktansatz
\begin{equation}
U(r, \varphi, z)
=
R(r)\, \Phi(\varphi)\, Z(z)
\label{bessel:equation:produktansatz}
\end{equation}
und setzen ihn in die Helmholtz-Gleichung ein.
Nach Division durch $U$ und Sortieren der Terme erhält man drei gewöhnliche Differentialgleichungen, je eine für $R(r)$, $\Phi(\varphi)$ und $Z(z)$.

Aus der $\varphi$-Gleichung folgt mit der Forderung nach $2\pi$-Periodizität, dass die Separationskonstante die Form $\nu^2$ mit ganzzahligem $\nu$ haben muss
% [TODO: Allgemeinere Argumentation bei nicht-ganzzahligem nu]
und $\Phi(\varphi) = e^{i\nu\varphi}$ als Lösung hat.
Die $z$-Gleichung führt auf Exponential- oder Trigonometrische Funktionen.

\subsection{Die Bessel'sche Differentialgleichung
\label{bessel:subsection:bessel-dgl}}
% Inhalt:
% - Aus der verbleibenden Radialgleichung wird durch Substitution x = kr
%   die "Standardform" der Bessel-DGL.
% - Endergebnis sauber als Satz oder hervorgehobene Gleichung.

Für den Radialteil $R(r)$ verbleibt nach der Separation die Gleichung
\begin{equation}
r^2 R''(r) + r R'(r) + (k^2 r^2 - \nu^2) R(r) = 0.
\label{bessel:equation:radial}
\end{equation}
Mit der Substitution $x = kr$ und $y(x) = R(x/k)$ geht diese Gleichung in die Standardform
\begin{equation}
\boxed{\;
x^2 y''(x) + x y'(x) + (x^2 - \nu^2) y(x) = 0\;
}
\label{bessel:equation:bessel-dgl}
\end{equation}
über.
Dies ist die \emph{Bessel'sche Differentialgleichung} der Ordnung $\nu$.
Sie ist von zweiter Ordnung, linear und homogen, und besitzt im Ursprung $x = 0$ eine reguläre Singularität.
Wegen dieser Singularität ist der einfache Potenzreihenansatz aus Kapitel~\ref{chapter:differentialgleichungen} nicht direkt anwendbar -- der nächste Abschnitt erweitert ihn entsprechend.

